\documentclass{new-aiaa}
\usepackage{amsmath}
\usepackage{graphicx}
\usepackage{booktabs}
\usepackage{longtable,array}
\usepackage{xcolor}
\newcommand{\gpc}[1]{\textcolor{teal}{\small\textsf{[#1]}}} % GP-class tag

\title{Geometric-Programming Sizing of a Single Oblique-Wing Trans-Medium Vehicle:\\ Model Formulation and Required CFD Inputs}
\author{Yiyang Zhang\footnote{Undergraduate Researcher, Harbin Institute of Technology.}}

\begin{document}
\maketitle

\begin{abstract}
We formulate the conceptual sizing of a torpedo-shaped hybrid aerial--underwater vehicle (HAUV)
carrying a single, centrally-pivoted oblique wing as a geometric program (GP). The wing is fully
deployed in aerial cruise, rotates during the dive, and lies parallel to the hull underwater, where
its overhang acts as a low-aspect-ratio lifting surface that trims the vehicle in depth and thereby
decouples weight from buoyancy. All analytical constraints are sourced from the marine-hydrodynamics
and GP-method literature (Renilson~\cite{renilson2018submarine}, Newman~\cite{newman2018marine},
Hoerner~\cite{hoerner1965fluid,hoerner1985fluid}, Brennen~\cite{brennen2014cavitation},
Myring~\cite{myring1976theoretical}, and the Hoburg--Abbeel GP framework~\cite{hoburg2014geometric,boyd2007tutorial});
residual non-analytic terms define a minimal CFD design of experiments (Sec.~\ref{sec:cfd}), to be
fitted into GP-compatible forms via the softmax/max-affine methods of Hoburg~et~al.~\cite{hoburg2016data}.
\end{abstract}

\section*{Nomenclature}
\noindent\begin{longtable}{@{}l p{10cm}@{}}
\multicolumn{2}{l}{\textbf{Geometry and configuration}}\\
$L$      & hull length (1.021~m from STL) \\
$d$      & hull diameter (0.100~m; fineness $L/d\approx 10$) \\
$S,\,b,\,c$ & wing planform area, span, chord \\
$\delta$ & wing overhang beyond hull when stowed (lifting length) \\
$\nabla$ & displaced hull volume; $S_{\text{wet}}$ wetted area \\
$A$      & wing aspect ratio $b^2/S$ \\
$\Lambda$& wing skew angle (deployed $\Lambda=0^\circ$, stowed $\Lambda=90^\circ$; $\Lambda=90-\text{WingXX}$) \\
$\alpha$ & angle of attack \\
\multicolumn{2}{l}{\textbf{Fluid properties}}\\
$\rho_a,\,\rho_w$ & density of air (1.225), water (1000~kg/m$^3$) \\
$\mu_a,\,\mu_w$ & dynamic viscosity \\
$p_v$    & vapour pressure of water ($\approx2.34$~kPa at $20^\circ$C)\cite{brennen2014cavitation} \\
$p_\infty$& hydrostatic pressure $p_{\text{atm}}+\rho_w g h$ \\
\multicolumn{2}{l}{\textbf{Forces, coefficients, and dimensionless groups}}\\
$C_L,\,C_D$ & lift, drag coefficient; $C_{D0}$ zero-lift drag \\
$C_f$    & skin-friction coefficient ($\text{Re}$-dependent) \\
$C_{Dp}$ & profile (section) drag coefficient \\
$C_m$    & pitching-moment coefficient; $C_l$ rolling-moment ($l$ = lower-case L, not numeral 1) \\
$C_p$    & pressure coefficient; $-C_{p,\min}$ suction peak \\
$\mathrm{Re}$& Reynolds number $\rho V L/\mu$ \\
$\sigma$ & cavitation number $(p_\infty-p_v)/(\tfrac12\rho_w V_w^2)$ \\
$(1+k)$ & form factor (Hoerner~\cite{hoerner1965fluid}): $1+1.5(d/L)^{1.5}+7(d/L)^3$, posynomial \\
$e$      & Oswald (span) efficiency (deployed); $e'$ stowed \\
$k$      & induced-drag factor (lift-polar $k$, not the form-factor $k$ --- these are distinct) \\
\multicolumn{2}{l}{\textbf{Mission and mass}}\\
$V_a,\,V_w$ & aerial (50~m/s), underwater cruise speed (swept 2--20~m/s) \\
$R_a,\,R_w$ & aerial, underwater cruise range \\
$B$      & buoyancy force $\rho_w g\nabla$ \\
$\xi$    & buoyancy--weight trim fraction $\xi=(W-B)/B$, swept \\
$m$      & total vehicle mass; $m_{\text{pay}}$ payload; $m_b$ battery; $m_s$ structure (hull + skin + fins) \\
$m_{a}$ & added mass (axial $m_{11}$; $m_{a}=k_a\,\rho_w\nabla$~\cite{newman2018marine}; for our fineness $L/d\approx10$: $k_a\approx0.03$, negligible) \\
$W$      & weight $mg$ \\
$e_b$    & battery specific energy (system-level Li-ion $\approx 100$--$180$~Wh/kg~\cite{rutherford2008southampton}; cell-level $\approx 190$--$300$~\cite{lin2026energies}) \\
$\eta$   & propulsive chain efficiency (motor $\approx0.85$ brushless DC~\cite{allen2000propulsion}; prop $\approx0.81$ open-water~\cite{allen2000propulsion}) \\
$a_{\text{peak}}$ & peak water-entry deceleration ($\le20$~g constraint)\\
$P$      & power (subscript $a$ air, $w$ water); $P_{\text{hotel}}$ hotel load \\
$t$      & time (per phase); $h$ operating depth \\
$t_s$    & skin (hull shell) thickness; $\sigma_y$ yield stress; $E$ Young's modulus \\
\end{longtable}

\section{Introduction}
Trans-medium vehicles that cruise in air and then continue submerged face a dynamic-pressure ratio of
order $130\times$ between the two media, which is why a fixed lifting surface cannot be shared across
both~\cite{long2025entire}. The present concept resolves this with a single wing on one central vertical pivot: deployed for
aerial lift ($\Lambda=0^\circ$), skewed flat against the hull for the dive and underwater phases
($\Lambda=90^\circ$). While GP-based conceptual sizing is well established for
aircraft~\cite{hoburg2014geometric,boyd2007tutorial,kirschen2018application}, no prior GP formulation
exists for an underwater or trans-medium vehicle. Constructing one, and identifying exactly which
coefficients must come from CFD, is the purpose of this model.

The GP conventions follow Hoburg and Abbeel~\cite{hoburg2014geometric} and
Boyd~et~al.~\cite{boyd2007tutorial}. The data-FITTING step that converts CFD outputs to GP-compatible
posynomials uses the softmax-affine and implicit-posynomial methods of
Hoburg~et~al.~\cite{hoburg2016data} and Kirschen~et~al.~\cite{kirschen2018application}.

\section{Geometric Programming Background}\label{sec:gp}
A \emph{monomial} is a function $h(\mathbf{u})=c\prod_{j}u_j^{a_j}$ with $c>0$ and
$a_j\in\mathbb{R}$~\cite{boyd2007tutorial}; a \emph{posynomial} is a sum of
monomials~\cite{hoburg2014geometric}. A GP minimizes a posynomial objective subject to
posynomial inequality constraints ($f_i(\mathbf{u})\le 1$) and monomial equality constraints.
Key modeling devices:

\begin{enumerate}
\item \textbf{Posynomial-equality relaxation.} A relation $y=f(\mathbf{u})$ with $f$ posynomial is
written $y\ge f(\mathbf{u})$; if the objective and all constraints involving $y$ are monotone in $y$,
the relaxation is tight at the optimum~\cite{hoburg2014geometric}.
\item \textbf{Maximizing by inversion.} Maximizing a quantity is equivalent to minimizing its monomial
inverse ($\min W^{-\alpha}$)~\cite{hoburg2014geometric,boyd2007tutorial}.
\item \textbf{Data-to-posynomial fitting.} Where no analytic form exists (e.g.\ CFD-generated polars),
a posynomial is fitted offline via softmax-affine / implicit-posynomial GP-compatible
regression~\cite{hoburg2016data}. This is the bridge between CFD and Layer~A.
\item \textbf{Signomial fallback.} Relations with a genuine subtraction (e.g.\ a force balance where a
lift term acts opposite to buoyancy) are handled as signomial programs (SP) or, when possible, made
GP-legal by sweeping a parameter (e.g.\ $\xi$) as a fixed constant per solve.
\end{enumerate}

\section{Mission and Configuration}
The mission has five phases: booster launch (from ship or aircraft), aerial cruise
($R_a=10$--$30$~km~at~50~m/s), unpowered dive, water entry, and underwater cruise
($R_w$). The booster is external and jettisoned — not sized here. The vehicle sizes for aerial
cruise, entry survival, and underwater cruise. Baseline geometry (Table~\ref{tab:geom}) is measured
from the STL export; the GP treats these as fixed inputs (``hybrid'' scheme).

\begin{table}[h]\centering
\caption{Baseline geometry measured from the STL export (metres).}\label{tab:geom}
\begin{tabular}{@{}llll@{}}\toprule
Quantity & Symbol & Value & Source \\\midrule
Hull length & $L$ & 1.021 & \texttt{NBody.stl} \\
Hull diameter & $d$ & 0.100 & \texttt{NBody.stl} \\
Wing chord & $c$ & $\approx0.14$ & \texttt{Wing90.stl} \\
Wing span (deployed) & $b$ & $\le0.72$ & constraint, Eq.~\eqref{eq:span} \\
Overhang / stow thickness & $\delta$ & 0.018 & \texttt{Wing*.stl} \\
Displaced volume & $\nabla$ & $\approx\!\pi d^2 L/4\approx8.0\times10^{-3}$~m$^3$ & CAD \\\bottomrule
\end{tabular}
\end{table}

A hard packaging constraint requires the stowed wing to fit along the hull with clearance:
\begin{equation}\label{eq:span}
b + 0.3\ \le\ L. \gpc{posynomial $\le$ monomial: GP-legal}
\end{equation}

\section{Geometric-Programming Model}\label{sec:gpmodel}

\subsection{Objective and requirements}
The objective is to \emph{maximize payload} for the required ranges:
\begin{equation}
\text{minimize}\quad m_{\text{pay}}^{-1},\qquad R_a\ge R_a^{\min},\ \ R_w\ge R_w^{\min}.
\gpc{monomial objective}
\end{equation}

\subsection{Aerial cruise: steady level flight}
With the wing deployed, lift balances weight, thrust balances drag~\cite{hoburg2014geometric}:
\begin{align}
m g &\le \tfrac12\rho_a V_a^2 S\,C_{L,a}, \gpc{monomial} \\
D_a &\ge \tfrac12\rho_a V_a^2 S\,C_{D,a}, \gpc{monomial (relaxed)} \\
P_a &\ge D_a V_a / \eta_a. \gpc{monomial}
\end{align}
These are the aircraft steady-flight relations of Ref.~\cite{hoburg2014geometric}, Eqs.~(15)--(18).

\subsection{Aerodynamic drag build-up}
The drag coefficient decomposes into induced, profile, and form
contributions~\cite{hoburg2014geometric,hoerner1965fluid}:
\begin{align}
C_{D,a} &\ge \frac{C_{L,a}^2}{\pi e A}
           + C_{Dp}(C_{L,a},\mathrm{Re},\Lambda)
           + C_{D0,a}, \gpc{posynomial}\\[4pt]
C_{D0,a} &= C_f\,(1+k)\,\frac{S_{\text{wet}}}{S},
           \qquad C_f = \frac{0.074}{\mathrm{Re}^{0.2}},
           \qquad \mathrm{Re} = \frac{\rho_a V_a c}{\mu}, \gpc{monomial}\\[4pt]
(1+k) &= 1 + 1.5\left(\frac{d}{L}\right)^{3/2}
        +     7.0\left(\frac{d}{L}\right)^{3}. \gpc{posynomial}\label{eq:formfactor}
\end{align}
Equation~\eqref{eq:formfactor} is Hoerner's~\cite{hoerner1965fluid} Eq.~(28) for streamlined bodies of
revolution. The term $1.5(d/L)^{3/2}$ captures the supervelocity increment (Eq.~26, $\Delta q_{\text{av}}/q
\propto(d/L)^{3/2}$); the $7(d/L)^3$ term is the separation/pressure component, negligible for
$d/L\lesssim0.2$. For our hull ($d/L\approx0.10$): $(1+k)\approx1.054$, i.e.\ a
\emph{5.4\%} viscous drag adder from thickness — a far smaller contribution than the $\sim$20\%
assumed in the project design notes, which conservatively included appendages. The SLA (service-life-allowance)
roughness/appendage increment is added as an empirical factor $f_{\text{SLA}}$, i.e.\
$(1+k)_{\text{eff}}=f_{\text{SLA}}(1+k)_{\text{naked}}$. The turbulent flat-plate law
$C_f=0.074\,\mathrm{Re}^{-0.2}$~\cite{hoburg2014geometric} is monomial and used in place of the ITTC
1957 log-law~\cite{myring1976theoretical,allen2000propulsion}; for $\mathrm{Re}\approx
2\times10^{7}$--$3\times10^{6}$ (water/air on $L\approx 1$~m) the power-law fit error is $<2\%$.\par
The profile-drag term $C_{Dp}(C_{L,a},\mathrm{Re},\Lambda)$ is \emph{not analytic} and must be supplied
as a fitted posynomial from CFD (Sec.~\ref{sec:cfd}), following the gpfit workflow of
Ref.~\cite{hoburg2016data}. The $\Lambda$-dependence of the Oswald factor $e(\Lambda)$ is also CFD-fitted;
at $\Lambda=0$ the analytic prior is $e\approx0.9$ from lifting-line theory~\cite{hoerner1985fluid}.

\subsection{Propulsion and endurance (electric)}
For a battery-electric vehicle the fuel-burn/Breguet model of Ref.~\cite{hoburg2014geometric} is
replaced by a constant-mass energy balance; endurance $t$ equals stored energy over
power~\cite{allen2000propulsion,rutherford2008southampton,lin2026energies}:
\begin{align}
E_b   &= e_b\, m_b,                  & R_a   &= V_a t_a,\quad R_w = V_w t_w, \gpc{monomial}\\
E_b   &\ge P_a t_a + P_w t_w + P_{\text{hotel}}t_{\text{total}}, \gpc{posynomial}\\
\eta_a&\le \eta_{\text{motor}}\,\eta_{\text{prop}}. \gpc{monomial}
\end{align}
For the aerial leg, $\eta_{\text{prop}}$ follows the actuator-disk model of
Ref.~\cite{hoburg2014geometric}, Eq.~(37), with $\eta_i\le 2/(1+\sqrt{1+T/(\frac12\rho V^2
A_{\text{prop}})})$. For the underwater leg, the optimal propeller efficiency is approximately
constant $\eta_{\text{prop,uw}}\approx0.81$~\cite{allen2000propulsion}. System-level Li-ion
specific energy is $\approx100$--$180$~Wh/kg (cell: $190$--$300$~Wh/kg)~\cite{lin2026energies};
REMUS field data~\cite{allen2000propulsion} confirm $1000$~Wh over $20$~h at $1.5$~m/s ($P
\approx50$~W). Hotel load $\approx 200$--$220$~W (survey-class AUV)~\cite{lin2026energies}; for
this one-off sprint design, factor it into the fixed-mass term or drop if payload power is
dominant.

\subsection{Mass build-up}
\begin{align}
m &\ge m_{\text{pay}} + m_s + m_b + m_{\text{prop}} + m_{\text{fixed}}, \gpc{posynomial}\\
m_{\text{prop}} &\le c_m\,P_{\max}. \gpc{monomial (power-law)}\\
m_s &\ge m_{\text{skin}} + m_{\text{fins}} + m_{\text{pivot}}. \gpc{posynomial}
\end{align}
If the hull is unstiffened (plausible at $100$~mm diameter), $m_{\text{skin}} = \rho_{\text{mat}}\,\pi
d\,L\,t_s$ with $t_s$ sized by the buckling gate (Sec.~\ref{sec:structure}). The power-law mass
coefficient $c_m$ (kg/kW) for a brushless DC motor + controller is $\approx 0.1$--$0.3$~kg/kW
(representative UAV/UUV values); verify against specific vendor data.

\subsection{Underwater cruise: buoyancy and wing trim}\label{sec:buoyancy}
Underwater the wing is stowed. Buoyancy supports most of the weight; the stowed overhang provides a
hydrodynamic trim force that closes the vertical balance. With $B=\rho_w g\nabla$ and the imbalance
fraction $\xi=(W-B)/B$ \emph{swept as a parameter}, the balance is monomial per solve:
\begin{align}
B   &= \rho_w g\, \nabla, \gpc{monomial}\\
m g &= B\,(1+\xi),        \gpc{monomial in $\xi$ (fixed per solve); SP if $\xi$ free}\\
L_{\text{tr}} &= \xi B = \tfrac12\rho_w V_w^2 S_{\text{stow}}\,C_{L,\text{uw}}. \gpc{monomial}
\end{align}
Sweeping $\xi$ from slightly negative (positively buoyant, wing trims down) to positive (heavy,
wing lifts — payload exceeds the buoyancy ceiling) traces a payload--trim-drag frontier. The trim
lift $L_{\text{tr}}$ induces drag $D_{\text{tr}}\ge k_{\text{uw}} L_{\text{tr}}^2/
(\tfrac12\rho_w V_w^2 S_{\text{stow}})$ that feeds the underwater power term.
The stowed wing's hydrodynamic behaviour is self-compensating: as the wing folds flat its aspect
ratio collapses $A\sim3.3 \to 0.1$ (Jones slender-body limit $\sim\pi A/2$~\cite{hoerner1985fluid}),
which reduces $C_{L\alpha}$ by roughly the same factor ($\sim22$) as the water/air $q$-ratio
($\sim130$) divided by the area reduction ($\sim3$): $130/(3\cdot22)\approx1.9$. This
near-cancellation — $\textit{not}$ an exact equality — is the analytic backbone of the
self-compensation thesis in the project design notes and is tested directly by the HYDRO CFD campaign.
The added-mass correction for vertical motion is negligible in the axial direction
($k_a\approx0.03$ at $L/d\approx10$~\cite{newman2018marine}) but may be $\sim$95\% of the
displaced mass laterally~\cite{newman2018marine}; on entry the effective mass $(m+m_a)$ replaces
$m$ in the impact constraint (Sec.~\ref{sec:entry}) and in any lateral-belly-flop transient (not
modelled here).

\subsection{Underwater hydrodynamic drag}
For the slender body of revolution,
\begin{align}
D_w      &= \tfrac12\rho_w V_w^2 S_{\text{wet}}\,C_{D,w}, \gpc{monomial}\\
C_{D,w}  &\ge C_f\,(1+k) + k_{\text{uw}}\,
            \frac{L_{\text{tr}}^2}{(\tfrac12\rho_w V_w^2)^2 S_{\text{stow}}^2}, \gpc{posynomial}
\end{align}
with $(1+k)$ from Eq.~\eqref{eq:formfactor} (same form factor, water Reynolds number). The trim-lift
induced-drag coefficient $k_{\text{uw}}$ comes from CFD (stowed wing, $\Lambda=90^\circ$). At
$d/L=0.10$, the minimum-drag thickness ratio $\approx 0.18$--$0.20$~\cite{hoerner1965fluid}; our
hull falls on the thinner (friction-dominated) side.

\subsection{Cavitation margin}\label{sec:cavitation}
The cavitation inception condition~\cite{brennen2014cavitation} is
\begin{equation}
-C_{p,\min}\ \le\ \sigma_i,
\qquad \sigma = \frac{p_{\text{atm}} + \rho_w g h - p_v}{\tfrac12\rho_w V_w^2},
\qquad p_v(20^\circ\text{C})\approx 2.34\ \text{kPa}.
\gpc{posynomial $\le$ monomial}
\end{equation}
At the surface ($h=0$, $V_w=20$~m/s): $\sigma\approx0.495$, consistent with
the project design notes~\S3. Amromin and Rozhdestvensky~\cite{amromin2022cavitation} show that at hull Reynolds
numbers $\mathrm{Re}\gtrsim10^7$ (our $\mathrm{Re}\approx2\times10^7$) the difference
$|\min C_p|-\sigma_i\to0$, i.e.\ the $\sigma_i=-C_{p,\min}$ rule is increasingly accurate. The
constraint is therefore conservative; the suction peak $-C_{p,\min}$ is the \emph{CFD deliverable}
(fitted vs $C_L$ or $\alpha$ from the HYDRO campaign). The GP enters it as
$\tfrac12\rho_w V_w^2(-C_{p,\min})+p_v\le p_{\text{atm}}+\rho_w g h$, which is posynomial in $V_w$
(with $-C_{p,\min}$ pre-fitted as a monomial/posynomial in $C_L$).

\subsection{Pressure-hull structure and buckling}\label{sec:structure}
For operating depth $h$, the thin-shell hoop stress (unstiffened tube) gives a monomial gauge
constraint~\cite{shinoka2024structural}:
\begin{equation}
\sigma_h = \frac{p_d\, d}{2 t_s} \le \sigma_y \ \Rightarrow\
p_d \le \frac{2\sigma_y t_s}{d},
\qquad p_d = \rho_w g h. \gpc{monomial}
\end{equation}
External-pressure buckling adds a second limit. The Bryant overall-elastic-buckling collapse
pressure~\cite{shinoka2024structural} is
\begin{equation}
p_{\text{cr}} = \frac{E\,t_s/a}{\,n^2-1+\lambda^2/2\ }
                \left[\frac{1}{(n^2+ \lambda^2)^2}\right] + \cdots
                \gpc{signomial (discrete $n$) / non-GP}
\end{equation}
where $\lambda=\pi a/L$, $a=d/2$ the radius, $n$ the integer buckling lobe number ($\ge2$).
The full Bryant form (not shown in full here) contains an additional $\lambda^4$ factor in the
numerator; this displayed sketch omits it — the equation is illustrative and will be replaced by a
fitted monomial in the GP implementation.
This is \emph{signomial} (posynomial denominator inverted) and has a discrete inner minimization over
$n$. For a thin unstiffened tube at our small diameter ($100$~mm), the recommended GP path is to
\emph{factor out} the $n$-minimization into an offline evaluation and fit a monomial
$p_{\text{cr}} \approx C\,E\,(t_s/a)^\alpha\,(a/L)^\beta$ over the $(t_s/a,\ L/a)$ range,
following the Windenburg--Trilling closed form. The operational safety factor is
$\mathrm{SF}_{\text{overall}}\ge2.0$~\cite{shinoka2024structural}; the constraint in the GP is
then $p_d \le p_{\text{cr}}/\mathrm{SF}$, which remains monomial after the fit.
Material: for a first pass use aluminium 6061-T6 ($E\approx69$~GPa, $\sigma_y\approx 250$~MPa,
$\rho\approx2700$~kg/m$^3$); carbon-fibre or HY-80 steel changes only the constants.

\subsection{Water-entry survival}\label{sec:entry}
The peak deceleration on water impact is modelled by a monomial fitted from the 2-D interFoam DOE:
\begin{equation}
a_{\text{peak}} = C\, V_{\text{entry}}^a\, d^b\, m_{\text{eff}}^c,
\qquad m_{\text{eff}} = m + m_a, \gpc{monomial fit}
\end{equation}
where $m_a=k_a\rho_w\nabla$ is the axial added mass
($k_a\approx0.03$~\cite{newman2018marine}; a $\sim$3\% correction). The analytic prior from
slamming theory is $a\approx2$ (force $\propto\frac12\rho V^2$ area), $b\approx2$ (diameter$^2$, from area $\propto d^2$),
$c\approx-1$ (Newton). All three exponents are \emph{fitted from CFD} — the prior serves only as a
sanity check. The hard constraint is $a_{\text{peak}}\le 20\,g$; entry angle enters as a secondary
parameter \emph{treated conservatively} (near-vertical worst case) unless the DOE returns strong
angular dependence. Song~et~al.~\cite{song2020waterentry} confirm the $V^2$ scaling for oblique
entry of high-speed projectiles. The GP constraint is monomial: $C V_{\text{entry}}^a d^b
m_{\text{eff}}^c \le a_{\max}$.

\subsection{Wing skew dynamics (residual roll — trajectory gate)}\label{sec:roll}
During the skew from $\Lambda=0^\circ$ to $\Lambda=90^\circ$, an asymmetric wake induces a roll
moment quantified by $C_l(\Lambda)$. This is the single highest-value CFD unknown — currently
unmeasured across a factor of $\sim20\times$ range. The roll-transient time constant is
$\tau\approx0.10$--$0.12$~s (the project design notes~\S4). The GP does \emph{not} optimize the skew trajectory
(that is Layer~B / Bayesian optimization); it only imposes the roll-authority feasibility gate:
\begin{equation}
|C_l(\Lambda)| \le C_{l,\max},
\qquad C_{l,\max} \triangleq \frac{\text{tail-fin authority}}{\tfrac12\rho V^2 S b}.
\gpc{monomial bound}
\end{equation}
$C_l(\Lambda)$ is supplied as a fitted posynomial vs $\Lambda$ from the AERO CFD campaign, and
$C_{l,\max}$ is computed from fin geometry: the cruciform ``VFin'' + ``HFin'' pair
($\sim$0.20~m span, $\sim$0.05~m chord) provide a roll-authority moment arm through a
lifting-line fin model. Feedforward tail-deflection scheduling further relaxes the gate (pre-program a
cancelling deflection per the known $C_l(\Lambda)$ function).

\subsection{Static stability and fin sizing}\label{sec:stability}
A trans-medium vehicle without a static-stability constraint could converge to an unflyable optimum.
Renilson~\cite{renilson2018submarine} (Ch.~3, Eq.~113, Table~3.8) defines dimensionless stability
indices in the vertical ($G_V$) and horizontal ($G_H$) planes:
\begin{align}
G_V &= 1 - \frac{M'_w\,(m' + Z'_q)}{M'_q\,Z'_w},
\qquad
G_H = 1 - \frac{N'_v\,(m' - Y'_v)}{N'_r\,Y'_v}\ \text{(by analogy)},\ \gpc{signomial (mixed derivative signs); $G_H$ inferred from vertical-plane form, not in Renilson OCR}\\[4pt]
\text{acceptable:}&\quad G_V\in[0.5,\,0.8],\quad G_H\in[0.2,\,0.4] \qquad\text{(Ray~et~al.~2008)}.
\end{align}
Primes denote non-dimensional stability derivatives divided by $\frac12\rho V L^2$ (e.g.\
$Z'_w = Z_w/(\frac12\rho V L^2)$). The derivatives themselves ($Z'_w, M'_w, Z'_q, M'_q$ for the
vertical plane; $Y'_v, N'_v, N'_r$ for the horizontal plane) can be approximated from fin geometry
via low-aspect-ratio lifting-surface theory~\cite{hoerner1985fluid}: for a cruciform fin set at the
tail, $Z'_w \propto (S_{\text{fin}}/S)(l_{\text{fin}}/\bar{c})$ — the tail volume
coefficient~\cite{hoerner1985fluid} (Ch.~XI),
\begin{equation}
V_H = \frac{S_{\text{fin}}}{S}\,\frac{l_{\text{fin}}}{\bar{c}},
\qquad
C_{L,\text{fin}} \propto V_H\,\alpha_{\text{fin}}. \gpc{monomial}
\end{equation}
For the GP, the $G_V$ and $G_H$ constraints enter as \emph{feasibility gates} with the derivatives
expressed as monomials in fin area and moment arm, using the analytic low-AR lift slope
$\partial C_L/\partial\alpha \approx \pi A_{\text{fin}}/2$~\cite{hoerner1985fluid} ($A_{\text{fin}}
\approx 2\cdot\text{span}_f/\text{chord}_f$). This is the single most important block missing from
the original sizing model. At our small scale ($L\approx1$~m), the cruciform fin pair
(span $\sim\!0.20$~m, chord $\sim\!0.05$~m) gives $A_{\text{fin}}\approx4$, yielding adequate
control authority; the stability gate validates this rather than leaving it unchecked.

\subsection{Propeller efficiency decomposition}\label{sec:propeller}
The GP currently uses a lumped propeller efficiency $\eta_{\text{prop}}\approx0.81$
\cite{allen2000propulsion}. Renilson~\cite{renilson2018submarine} (Ch.~5, Eq.~5.4) decomposes the
quasi-propulsive coefficient (QPC) into physically distinct factors:
\begin{equation}
\text{QPC} = \eta_H \cdot \eta_O \cdot \eta_R,
\qquad \eta_H = \frac{1-t}{1-w}, \gpc{monomial}
\end{equation}
where $t$ = thrust deduction fraction, $w$ = Taylor wake fraction (hull-form dependent,
$\eta_H\approx1.05$ for streamlined bodies~\cite{renilson2018submarine}), $\eta_O$ = open-water
propeller efficiency (function of advance ratio $J$ and loading $K_T/J^2$), and $\eta_R$ = relative
rotative efficiency $\approx1.05$ for single-propeller configurations
($D_{\text{prop}}/D_{\text{hull}}\approx0.4$--$0.7$). The open-water propeller coefficients are
canonical monomials~\cite{newman2018marine} (Ch.~2):
\begin{equation}
K_T = \frac{T}{\rho n^2 D^4},\quad
K_Q = \frac{Q}{\rho n^2 D^5},\quad
J   = \frac{U_A}{n D},\quad
\eta_O = \frac{J}{2\pi}\frac{K_T}{K_Q}. \gpc{monomial}
\end{equation}
If propeller diameter $D$ and RPM $n$ become GP variables, $K_T$ and $K_Q$ are monomials that link
thrust/torque to $D$ and $n$; the design-space constraint $K_T/J^2 \le (K_T/J^2)_{\max}$ prevents
overloading. For the fixed-geometry (``hybrid'') baseline, treat $\eta_O$, $\eta_H$, $\eta_R$ as
constants and revert to $\text{QPC}\approx0.81$. Common AUV system-level values:
motor $\eta_m\approx0.85$ (brushless DC), QPC $\approx0.8$--$0.9$~\cite{renilson2018submarine},
yielding the chain $\eta_{\text{prop}} = \eta_m\eta_{\text{QPC}} \approx 0.7$--$0.75$.

\subsection{Summary of the GP}
The full GP (Layer~A) is:
\begin{equation}\label{eq:gpsummary}
\begin{array}{ll}
\text{minimize}    & m_{\text{pay}}^{-1} \\
\text{subject to} & \text{Eqs. (steady flight, drag, energy, mass, buoyancy/$\xi$, hydro drag,} \\
                  & \hspace{2pt}\text{cavitation, hoop stress, buckling [fitted], entry, roll gate,} \\
	                  & \hspace{2pt}\text{stability $G_V,G_H$, propeller QPC)} \\
\text{with}        & \text{$\xi$ swept; $\Lambda=\text{fixed}$ (solved at each $\Lambda$)} \\
                  & C_{Dp}(\Lambda),\ k(\Lambda),\ C_{L,\text{uw}},\ k_{\text{uw}},\ -C_{p,\min},\ C_l(\Lambda),\ a_{\text{peak}} \text{ from CFD}.
\end{array}
\end{equation}

\subsection{Sizing constants and solved design point}
The Layer-A GP is implemented in \texttt{gpkit}~\cite{hoburg2014geometric} (\texttt{uauv\_gp\_model.py});
Table~\ref{tab:const} lists the sizing constants and Table~\ref{tab:soln} the solved design point at
buoyancy imbalance $\xi=0.15$. Water-entry is excluded from the GP and applied as a separate
post-hoc survival check. The model is a pure GP with $\xi$ swept; freeing $\xi$ makes
$mg=B(1+\xi)$ signomial (solved by \texttt{localsolve}).
\begin{table}[h]\centering
\caption{Sizing constants (SI). Geometry from the CAD; coefficients from the CFD fits.}\label{tab:const}
\small
\begin{tabular}{@{}llll@{}}\toprule
Symbol & Value & Symbol & Value \\\midrule
$S$ (wing/ref area) & $0.075$~m$^2$ & $\nabla$ (displaced vol.) & $0.0052$~m$^3$ \\
$d,\ L$ (hull) & $0.10,\ 1.021$~m & $S_{\text{wet}}$ & $0.42$~m$^2$ \\
$V_a,\ V_w$ & $50,\ \ge10$~m/s & $R_a,\ R_w$ & $20,\ 2$~km \\
$\eta_a,\ \eta_w$ & $0.70,\ 0.81$ & $e_b$ (Li-ion) & $0.54$~MJ/kg \\
$C_{D0,w},\ k_w$ (CFD) & $0.033,\ 21.3$ & $C_{L,\max}$ (CFD) & $0.60$ \\
$-C_{p,\min}$ (CFD) & $1.16$ & $k'_a$ polar (CFD) & $0.55$ \\
$\sigma_y,\ E$ (Al) & $250,\ 69{\times}10^3$~MPa & depth $h$ & $10$~m \\\bottomrule
\end{tabular}
\end{table}
\begin{table}[h]\centering
\caption{Solved Layer-A design point ($\xi=0.15$, maximise payload).}\label{tab:soln}
\small
\begin{tabular}{@{}llll@{}}\toprule
Quantity & Value & Quantity & Value \\\midrule
payload $m_{\text{pay}}$ & $2.96$~kg & vehicle mass $m$ & $5.98$~kg \\
battery $m_b$ & $1.44$~kg & structure $m_s$ & $1.28$~kg \\
aerial $C_{L,a}$ & $0.51$ & aerial $C_{D,a}$ & $0.14$ \\
aerial power $P_a$ & $1.18$~kW & underwater power $P_w$ & $1.54$~kW \\
hull wall $t_s$ & $1.08$~mm & (govern: hoop stress) & \\\bottomrule
\end{tabular}
\end{table}
The dominant sensitivities are $\partial(m_{\text{pay}}^{-1})/\partial\nabla=-1.43$ (displaced volume
is the strongest payload lever), $\partial/\partial e_b=-0.49$ (battery energy density), and
$\partial/\partial\eta_a=-0.30$; the design is buoyancy-limited, and the feasible $\xi$ is bounded
above by the aerial stall ceiling $C_{L,a}\le C_{L,\max}$.

\section{Required CFD Inputs}\label{sec:cfd}
Every non-analytic term above is a CFD deliverable. These, and \emph{only} these, define the design
of experiments; all other quantities come from the analytic sources in Table~\ref{tab:cfd}.
\begin{table}[h]\centering
\caption{CFD coefficients, their GP constraints, and the DOE independent variables.}\label{tab:cfd}
\begin{tabular}{@{}llll@{}}\toprule
CFD output & GP constraint closed & Campaign & DOE sweep \\\midrule
$C_{D0}(\Lambda),\,k(\Lambda)$ & aerial drag polar & AERO & $\alpha\times\Lambda$ (5 skews: 0/30/45/60/90$^\circ$) \\
$C_{Dp}(C_L,\mathrm{Re},\Lambda)$ & profile drag posynomial & AERO & $\alpha\times\Lambda$ \\
$C_{L\alpha}(\Lambda)$ & lift/weight balance & AERO & $\alpha\times\Lambda$ \\
$C_l(\Lambda)$ & roll-authority gate & AERO & $\Lambda$ sweep (wing-only force budget) \\
\addlinespace
$C_{D0,w},\,k_w$ & underwater drag polar & HYDRO & $\alpha\times\Lambda + V_w$ sweep at deployed \& stowed \\
$C_{L,\text{uw}}(\alpha),\,k_{\text{uw}}$ & $\xi$-trim vertical balance & HYDRO & $\alpha$ at $\Lambda=90^\circ$ (stowed wing) \\
$-C_{p,\min}(C_L)$ & cavitation $\sigma_i\ge -C_{p,\min}$ & HYDRO & $\alpha\times V_w\times h$ \\
\addlinespace
$a_{\text{peak}}=C V_{\text{entry}}^a d^b m^c$ & impact $\le20\,g$ & ENTRY & $V_{\text{entry}}\times d\times m_{\text{eff}}$ (2-D) \\\bottomrule
\end{tabular}
\end{table}
The AERO campaign uses \texttt{simpleFoam} (steady, k-$\omega$ SST, incompressible, air 50~m/s) with
\emph{10} skew meshes (CAD Wing0--Wing90 in $10^\circ$ steps) and sweeps $\alpha\in[-4^\circ,14^\circ]$
by rotating the inlet velocity vector — \emph{no remeshing}. The HYDRO campaign uses \texttt{simpleFoam}
with water properties and adds a $V_w$ sweep (2, 5, 10, 15, 20~m/s) at deployed and stowed configurations,
plus depth sweeps ($h=0,2,5,10$~m) for the cavitation margin. The ENTRY campaign uses
\texttt{interFoam} VOF (2-D, one-cell-thick, cavitation OFF baseline) for a vertical slender-ogive nose
at entry velocities (30--70~m/s), angles ($-15^\circ$ to $-75^\circ$), and masses (3--8~kg). Each
fit uses the gpfit protocol of Hoburg~et~al.~\cite{hoburg2016data}: scale the data to
dimensionless form, fit the target expression, and report RMS error; return \texttt{NaN} on
divergence, and cache every evaluation keyed by the design vector (CSV / SQLite).

\subsection*{Fitted CFD coefficients (completed campaigns)}
The AERO and HYDRO campaigns are complete (100 and $\sim$190 converged runs); Table~\ref{tab:cfdfit}
lists the fitted values (from \texttt{cfd/database/*.csv} via \texttt{cfd/fit\_gp.py}; full report in
\texttt{gp\_build/fitted\_coefficients.md}). The stowed induced-drag factor cross-validates between
air and water to $7.1\%$ ($k=22.9$ vs $k_w=21.3$, Buckingham-$\Pi$), and the lift slope obeys
$C_{L\alpha}\propto(\cos\Lambda)^{1.105}$ (rms$_{\log}=0.011$), confirming the sweep-independence prior.
\begin{table}[h]\centering
\caption{CFD-fitted GP coefficients (2026-08-03). Convention $\Lambda=90^\circ-\text{Wing}XX$.}\label{tab:cfdfit}
\begin{tabular}{@{}lll@{}}\toprule
Coefficient & Fitted value & Fit quality \\\midrule
Deployed cruise point ($\Lambda{=}0$) & $C_L{=}0.365,\ C_D{=}0.0670$ & measured \\
Deployed induced factor $k_a$ & $0.638$ & $R^2{=}0.994$ \\
Deployed lift slope $C_{L\alpha,a}$ & $0.835$/rad\,$^\dagger$ & $R^2{=}1.000$ \\
Lift-slope schedule $C_{L\alpha}(\Lambda)$ & $0.847\,(\cos\Lambda)^{1.105}$ & rms$_{\log}{=}0.011$ \\
Stowed polar ($\Lambda{=}90$) & $C_{D0,w}{=}0.0332,\ k_w{=}21.3$ & $R^2{=}0.988$ \\
Underwater $C_{D0,w}(\mathrm{Re})$ & $0.185\,\mathrm{Re}^{-0.112}$ & rms$_{\log}{=}0.016$ \\
Stowed trim slope $C_{L\alpha,\text{uw}}$ & $0.053$/rad & $R^2{=}0.966$ \\
Cavitation peak $-C_{p,\min}$ & $1.132\,(1{+}C_L)^{-0.078}$ (max $1.16$) & rms$_{\log}{=}0.032$ \\
Entry $a_{\text{peak}}$ & pending (DOE running) & --- \\\bottomrule
\end{tabular}
\end{table}

\noindent$^\dagger$ \textbf{Deployed-wing absolute efficiency is not yet grid-converged.} The
$57$k-cell sweep mesh under-resolves the wing: $C_{L\alpha,a}=0.835$/rad is $\sim\!5\times$ below
thin-airfoil theory ($\approx4.8$ for $AR\!\approx\!6.7$), and $k_a=0.638$ implies an Oswald
efficiency $e=1/(\pi k_a AR)\approx0.07$ versus the expected $0.8$--$0.9$. The mesh-\emph{converged}
outputs are the \emph{dimensionless scalings} --- $C_{L\alpha}\!\propto\!(\cos\Lambda)^{1.11}$,
$k\!\propto\!(\cos\Lambda)^{-1.35}$, $C_{D0,w}\!\propto\!Re^{-0.11}$ --- the cruise $C_L$, the stowed
polar (genuine $AR\!\to\!0$ collapse), the cavitation gate, and the $7\%$ air/water cross-validation.
The absolute deployed $C_{L\alpha}$ and $e$ await either a grid-convergence mesh plus a wider
$\alpha$-sweep, or a section-level NACA-3612 polar for $C_{Dp}(C_L)$ (as in the authors' prior
work). The narrow measured range ($C_L\!\in\![0.30,0.57]$, upper branch only) also precludes a
full drag-bucket softmax-affine posynomial here. See \texttt{gp\_build/verify\_gpfit.md}.

\section{Implementation and Validation Plan}
\begin{enumerate}
\item \textbf{gpkit prototype} with placeholder coefficients (implemented in \texttt{gp\_layerA.py}):
  solve at each $\xi$ and each $\Lambda$, verifying the solver converges and the
  sensitivity-tornado (free dual variables) identifies which constraints bind.
\item \textbf{CFD DOE} per Table~\ref{tab:cfd}: produce the fitted posynomials for every non-analytic
  term, with a grid-convergence study (1--8~Mcells, Richardson-extrapolation asymptotic range) at
  the nominal design point.
\item \textbf{GF \& SP pass}: substitute the CFD fits into the GP, re-solve, and produce the
  payload--mass frontier, the $\delta$-trade curve ($\delta$ vs trim drag), and the $\xi$-sweep
  self-compensation crossover chart.
\item \textbf{Layer~B trajectory optimization}: with the sized geometry and fitted coefficients,
  Bayesian optimization over the skew schedule ($\Lambda$-onset altitude, skew rate
  $d\Lambda/dt$, coast gap, entry/exit angles) using the 3-DOF trajectory sim.
\item \textbf{Citation reconciliation}: all entries in \texttt{refs/uauv.bib} are pinned to exact
  PDF page numbers (verified from the OCR corpus); Zotero keys follow the
  \texttt{auth:lower[year][firsttitleword:lower]} convention.
\end{enumerate}

\bibliographystyle{new-aiaa}
\bibliography{refs/uauv}
\end{document}
